\section{Verification}\label{sec:verification}

\paragraph{Exact core.}
All quantities are computed by the module \texttt{ws\_core.py} from the primitives of Definition~\ref{def:system}: integers and bitmasks for the finite objects, \texttt{fractions.Fraction} for the continuous fibre. The belief game of Definition~\ref{def:game} is implemented once, generically, as a map from nodes to option families; the kernel is the greatest fixed point on the reachable node graph, obtained by iterating $\Phi$ from the set of all reachable nodes until stabilisation (Lemma~\ref{lem:basic}(a)), so no horizon is chosen by hand. The design census uses a second, bitmask implementation of the free-protocol game on subsets of the eight viable states; the two implementations agree on the four classes. The ledger and private no-repeat kernels are computed a third time from the closed forms of Theorem~\ref{thm:closed}, by code that never calls the game engine, and agree with it on all $511$ beliefs of safe states for both observation maps.

\paragraph{What is checked.}
The verifier \texttt{paper11c\_worked\_systems\_audit\_v5\_verification.py} recomputes everything from the core (about $70$ seconds on one processor) and evaluates \NCHECKS{} named checks, each a Boolean statement about the recomputed objects: the sizes and set identities of Table~\ref{tab:kernels}; the strictness witnesses; the stabilisation horizons; the closed-form agreements; the stationary-policy unions; the decentralised factoring (the factored computation against the direct orbit computation), the $4\times4$ product, the authority census and the behavioural quotient; the design census with the closed-form count \eqref{eq:count}, the sixteen coarsest designs, the number of distinct kernels and the characterisation of the minimum; the four-state census, the inclusion--exclusion count and the Helly numbers; the regime kernels; the timing grid with the three boundary cells and the $21$ failing cells of the floored strict form, and the max-plus recursion against path enumeration; the benchmark crossover, margin law and radius; the duality values and the nonconvex gap; the drift identities; and the scaling table for caps $4$ and $5$. A second group of checks reads the compiled manuscript and the supplement and confirms that every number quoted in them matches the recomputed value. The verifier exits with a nonzero status on any failure; its complete output is reproduced in the supplement (Section~S10).

\paragraph{Mechanised layer.}
The monotonicity theorem is formalised in the repository's Lean~4 development, modules \texttt{Formalizations/WS\_WorkedSystems.lean} and \texttt{WS\_WorkedSystems\_v2.lean}, on top of the finite-system framework of the obstruction-calculus module; the theorems \texttt{Wk\_mono\_obs} (observation axis), \texttt{master\_action\_axis} (protocol axis), \texttt{master\_memory\_axis} and the canonical \texttt{master\_monotone} state the horizon-indexed inclusions of Theorem~\ref{thm:mono}. The verifier checks that these declarations are present in the modules shipped with the article; the build itself is documented in the repository.

\paragraph{Reproducibility.}
The article, the supplement, the figures and all tables are generated from \texttt{results.json}, which is written by \texttt{run\_all.py}; no number in the text is typed by hand. The scripts, the Lean modules and the verifier output are part of the submission package and are available at \url{https://github.com/MIKEAA2020/general-sustainability}.
